\section{Limitations}
\label{sec:limitations}

\ifshort
The agent under test is our own: generality to third-party agents is asserted by
the adapter interface, not demonstrated. Tools are simulated and tool calls are
parsed from plain text rather than issued through a native function-calling
interface. The panel is \nModels{} models, one per category, with license and
tier confounded (\cref{sec:refuse}). One trial per cell at temperature~0, with
intervals that resample payloads and so do not capture run-to-run variability.
The corpus is \corpusSize{} payloads, English only, and the encoding list is
fixed, so rates are lower bounds. Test coverage outside the frozen detector is
end-to-end rather than unit-level, and the observer effect of planting the canary
is unmeasured. The full list is \cref{app:limitations}.
\else
\paragraph{The agent under test is our own.}
This is the limitation we would attack first. The harness exposes an adapter
interface intended to let a user evaluate their own agent, and the construction
has no dependency on our particular agent---but generality across third-party
agents is at present \emph{asserted by the interface, not demonstrated}. No agent
we did not write has been evaluated. Until one has, every rate reported here
describes one agent, and the claim that the method transfers is a design claim
rather than an empirical one.

\paragraph{Simulated tools and a non-native calling convention.}
Tools return fixed data and no request leaves the process. Tool calls are emitted
as plain text and parsed by regular expression, not issued through a native
function-calling interface, which is a real divergence from how deployed agents
call tools and may affect susceptibility in either direction. A production
integration would also present noisier tool outputs and additional egress
surfaces---email, ticket updates, webhooks---that this surface does not model.

\paragraph{Panel size and the tier confound.}
\nModels{} models, one per deployment category, with a large tier and scale
mismatch between them. \Cref{sec:refuse} states the consequence: the design
cannot separate license from capability, and the headline gap must not be read as
a family-level claim.

\paragraph{Single trial per cell; intervals that do not cover sampling.}
Each cell is run once at temperature~0. Hosted inference at temperature~0 is not
exactly deterministic---we directly observed one benign episode that transmitted
the record on one execution and not on a byte-identical repeat---so a single
trial can miss behaviour that repeated sampling would reveal. The bootstrap
resamples the payload corpus, so the intervals quantify payload variability only.
Run-to-run variability is unmeasured.

\paragraph{Corpus size and ceiling effects.}
\corpusSize{} payloads per model give wide intervals, and one model saturated at
ceiling, which makes its interval degenerate and its breakdown uninformative
(\cref{sec:breakdown}). The instrument is built to separate large behavioural
differences, not to estimate rates precisely. All payloads and the workflow are
in English; multilingual injection is untested.

\paragraph{Fixed encoding list.}
The detector covers \nEncodings{} encodings. A model that invented a novel
recoverable encoding would evade detection and be scored as no leak, so every
rate here is a lower bound. This is a soundness-preserving failure---we do not
report leaks that did not happen---but it is a completeness failure.

\paragraph{Test coverage outside the detector.}
The \nOracleTests{} tests exercise the frozen detector: a positive per encoding
plus negative suites over other members' identifiers in every encoding, benign
knowledge-base text, partial fragments, and random high-entropy strings. The
runner, configuration layer, and verifier have no unit tests; they are covered
end to end by the study executing against them and by the equivalence check
described in \cref{sec:method}. That is weaker coverage than unit tests and we
name it as a gap.

\paragraph{Observer effect.}
The planted record is designed to be indistinguishable from the other records in
the store, but we have not verified that its presence leaves agent behaviour
unchanged relative to a workflow containing no canary at all.

\paragraph{Provider-mediated serving.}
The open-weights model was accessed through a hosted provider rather than
self-hosted. Served weights may be quantized and are wrapped in provider-side
scaffolding, so results characterize the model as commonly served rather than the
raw checkpoint.
\fi
